\documentclass[11pt,a4paper]{article}
% =====================================================================
%  Shared preamble for the Part V lecture notes (lectures 19-22).
%
%  These four lectures sit outside Geron, so the notes ARE the primary
%  source and are examined on the same terms as the textbook parts.
%  Everything here is chosen to make that readable on paper: one column,
%  generous measure, and the same six-colour palette as the slides so a
%  student moving between the two is not re-learning what red means.
%
%  Build with pdflatex (twice, for the toc and the refs):
%      cd notes && latexmk -pdf lecture-19.tex
%  or  make -C notes
% =====================================================================
\usepackage[T1]{fontenc}
\usepackage[utf8]{inputenc}
\usepackage{newtxtext,newtxmath}      % Times-like: prints well, reads well
\usepackage[scaled=0.86]{inconsolata} % code, at a size that sits beside Times
\usepackage{microtype}
% newtx defines \Bbbk, \openbox, \proof and \endproof; amssymb and amsthm
% then redefine them and abort. Freeing the four names before those two
% packages load is the standard fix, and it costs nothing here: we use
% amsthm only for \newtheorem.
\let\Bbbk\relax
\usepackage{amsmath,amssymb}
\let\openbox\relax
\let\proof\relax
\let\endproof\relax
\usepackage{amsthm}
\usepackage{booktabs}
\usepackage{enumitem}
\usepackage{xcolor}
\usepackage{tcolorbox}
\usepackage{titlesec}
\usepackage{graphicx}
\usepackage[hidelinks]{hyperref}
\tcbuselibrary{skins,breakable}

% ---- the course palette, from assets/css/custom.css -------------------
\definecolor{aimlprimary}{HTML}{0B3D62}   % deep blue  - structure
\definecolor{aimlaccent} {HTML}{C0392B}   % brick red  - failures
\definecolor{aimlsuccess}{HTML}{14663A}   % green      - repairs
\definecolor{aimlmath}   {HTML}{6C3483}   % purple     - the derivation
\definecolor{aimlmuted}  {HTML}{4B5563}
\definecolor{aimlrule}   {HTML}{B0BCC7}
\definecolor{aimlsoft}   {HTML}{F4F7F9}

\usepackage[a4paper,margin=32mm,top=30mm,bottom=30mm]{geometry}
\setlength{\parskip}{0.45\baselineskip}
\setlength{\parindent}{0pt}
\linespread{1.06}

% ---- headings --------------------------------------------------------
\titleformat{\section}{\Large\bfseries\color{aimlprimary}}{\thesection}{0.7em}{}
\titleformat{\subsection}{\large\bfseries\color{aimlprimary}}{\thesubsection}{0.7em}{}
\titleformat{\subsubsection}{\bfseries\color{aimlmuted}}{}{0em}{}
\titlespacing*{\section}{0pt}{1.6\baselineskip}{0.5\baselineskip}
\titlespacing*{\subsection}{0pt}{1.2\baselineskip}{0.35\baselineskip}

% ---- boxes -----------------------------------------------------------
% One box per role, and the role is the colour. A student who has seen the
% slides already knows what each one means before reading a word of it.
\newtcolorbox{keybox}[1][]{
  breakable, enhanced, colback=aimlsoft, colframe=aimlprimary,
  boxrule=0.4pt, left=8pt, right=8pt, top=6pt, bottom=6pt,
  fonttitle=\bfseries\small, coltitle=white, colbacktitle=aimlprimary,
  attach boxed title to top left={yshift=-2mm, xshift=4mm},
  boxed title style={boxrule=0pt, arc=1pt}, #1}

\newtcolorbox{failbox}[1][]{
  breakable, enhanced, colback=white, colframe=aimlaccent,
  boxrule=0.9pt, left=8pt, right=8pt, top=6pt, bottom=6pt,
  fonttitle=\bfseries\small, coltitle=white, colbacktitle=aimlaccent,
  attach boxed title to top left={yshift=-2mm, xshift=4mm},
  boxed title style={boxrule=0pt, arc=1pt}, #1}

\newtcolorbox{derivbox}[1][]{
  breakable, enhanced, colback=white, colframe=aimlmath,
  boxrule=0.6pt, left=8pt, right=8pt, top=6pt, bottom=6pt,
  fonttitle=\bfseries\small, coltitle=white, colbacktitle=aimlmath,
  attach boxed title to top left={yshift=-2mm, xshift=4mm},
  boxed title style={boxrule=0pt, arc=1pt}, #1}

% An exercise. Deliberately the same shape as the notebook's `try` field:
% one change, and what should happen to the output.
\newtcolorbox{trybox}[1][]{
  breakable, enhanced, colback=white, colframe=aimlsuccess,
  boxrule=0.5pt, left=8pt, right=8pt, top=6pt, bottom=6pt,
  fonttitle=\bfseries\small, coltitle=white, colbacktitle=aimlsuccess,
  attach boxed title to top left={yshift=-2mm, xshift=4mm},
  boxed title style={boxrule=0pt, arc=1pt}, #1}

% ---- small semantics -------------------------------------------------
\newcommand{\scope}[1]{\textcolor{aimlmuted}{\small #1}}
\newcommand{\term}[1]{\textbf{#1}}
\newcommand{\code}[1]{\texttt{\small #1}}
\newcommand{\lecture}[1]{Lecture~#1}

\newtheorem{definition}{Definition}[section]
\newtheorem{proposition}[definition]{Proposition}

% ---- title block -----------------------------------------------------
% \notestitle{19}{Information retrieval: the lexical foundation}{SciFact (BEIR)}
\newcommand{\notestitle}[3]{%
  \thispagestyle{empty}
  \begin{flushleft}
    {\color{aimlmuted}\small\sffamily
     APPLICATIONS OF MACHINE LEARNING \quad$\cdot$\quad
     BSc Mathematics of Artificial Intelligence}\\[2mm]
    {\color{aimlrule}\rule{\textwidth}{0.8pt}}\\[4mm]
    {\color{aimlmuted}\large Lecture #1 \quad$\cdot$\quad Part V \quad$\cdot$\quad
     Extended notes}\\[3mm]
    {\Huge\bfseries\color{aimlprimary}\baselineskip=1.15\baselineskip #2\par}\vspace{4mm}
    {\color{aimlmuted}\large Dataset: #3}\\[3mm]
    {\color{aimlrule}\rule{\textwidth}{0.8pt}}
  \end{flushleft}
  \vspace{2mm}
  \begin{keybox}[title={Why these notes exist}]
    Lectures 19--22 sit outside G\'eron, the course's primary text. For those
    four lectures \emph{these notes are the primary source}, and they are
    examined on exactly the same terms as the textbook parts. Everything in
    the slide deck for this lecture is here, with the arguments written out at
    length rather than compressed onto a slide; the numbers are the ones the
    accompanying notebook prints, and every one of them is a measurement on
    the corpus named above rather than a figure quoted from a paper.
  \end{keybox}
  \vspace{1mm}
  {\small\tableofcontents}
  \vspace{2mm}
  \noindent{\color{aimlrule}\rule{\textwidth}{0.4pt}}
  \vspace{1mm}
}


\begin{document}
\notestitlefor{2}{I}{The end-to-end project}{California housing}{Chapter 2, with \S4.1 brought forward}

\section{Where we are}

Lecture 1 looked at the data and split it. Five findings carry into today, and
three of them say the same thing.

\begin{center}
\small
\begin{tabular}{p{0.42\textwidth}p{0.46\textwidth}}
\toprule
\textbf{What we found} & \textbf{Why it matters today} \\
\midrule
One categorical column, \code{ocean\_proximity}, one level with $n=5$ &
it has to be encoded, and that level will be absent from some folds \\
207 missing values in \code{total\_bedrooms} &
something must fill them, and that something is \emph{learned from data} \\
Features on wildly different scales &
they have to be scaled, and scaling is \emph{learned from data} \\
A target capped at \$500{,}000 &
4.8\% of rows carry a label that is not the answer \\
16{,}512 training / 4{,}128 test, stratified on income &
the test rows are sealed until the last fifteen minutes \\
\bottomrule
\end{tabular}
\end{center}

Three of those five say \emph{learned from data}, and anything learned from
data can be learned from the wrong data. That is the subject of
\S\ref{sec:generalisation}.

\section{Derivation: least squares and the normal equation}

You called \code{LinearRegression().fit()}. What did it compute, and when does
that computation fail? The second half of the question is not idle: its answer
is exactly the warning that ended Lecture 1 about engineering features as
weighted sums of existing ones.

\subsection{The model}

\[ \hat{y} \;=\; \theta_0 + \theta_1 x_1 + \theta_2 x_2 + \dots + \theta_n x_n \]

with $\hat y$ the prediction, $n$ the number of features, $x_i$ the $i$-th
feature and $\theta_j$ the $j$-th parameter. $\theta_0$ is the \term{bias} or
\term{intercept}. Writing $\mathbf{x} = (x_0, \dots, x_n)$ with $x_0 \equiv 1$
folds the bias into the dot product:
\[ \hat{y} \;=\; h_{\boldsymbol\theta}(\mathbf{x}) \;=\;
   \boldsymbol\theta \cdot \mathbf{x} \;=\;
   \boldsymbol\theta^{\intercal}\mathbf{x}. \]

\subsection{What ``fit'' means}

Training means setting $\boldsymbol\theta$ so the model best fits the training
set, which requires a measure of how badly it fits:
\[ \mathrm{MSE}(\mathbf{X}, y, h_{\boldsymbol\theta}) \;=\;
   \frac{1}{m}\sum_{i=1}^{m}
   \left(\boldsymbol\theta^{\intercal}\mathbf{x}^{(i)} - y^{(i)}\right)^{2}. \]

We minimise MSE rather than RMSE because $\sqrt{\cdot}$ is strictly increasing
on $[0,\infty)$, so
\[ \arg\min_{\boldsymbol\theta} \sqrt{f(\boldsymbol\theta)}
   \;=\; \arg\min_{\boldsymbol\theta} f(\boldsymbol\theta): \]
same minimiser, simpler algebra. This is the first instance of a pattern worth
naming now, because it recurs for the rest of the course: \emph{learning
algorithms routinely optimise a different function during training than the one
used to evaluate the result} --- because it is easier to optimise, or because
it carries terms needed only during training.

\subsection{The normal equation}

Stack the instances as rows of $\mathbf{X}$ and the labels into $y$:
\[ \mathrm{MSE}(\boldsymbol\theta) \;=\;
   \frac{1}{m}\,\lVert \mathbf{X}\boldsymbol\theta - y \rVert_2^{\,2}. \]
We are minimising a squared Euclidean norm. That is a \emph{geometric}
statement, and it is the whole content of this section.

\begin{derivbox}[title={Differentiate, and check the second order}]
\[ \nabla_{\boldsymbol\theta}\,\mathrm{MSE}(\boldsymbol\theta) \;=\;
   \frac{2}{m}\,\mathbf{X}^{\intercal}\!
   \left(\mathbf{X}\boldsymbol\theta - y\right), \]
so at a stationary point $\hat{\boldsymbol\theta}$,
\[ \mathbf{X}^{\intercal}\!\left(\mathbf{X}\hat{\boldsymbol\theta} - y\right)
   \;=\; \mathbf{0}. \]
Is it a minimum? The Hessian is
$\nabla^{2}_{\boldsymbol\theta}\,\mathrm{MSE} =
 \frac{2}{m}\,\mathbf{X}^{\intercal}\mathbf{X}$,
and for any $\mathbf{v}$,
\[ \mathbf{v}^{\intercal}\mathbf{X}^{\intercal}\mathbf{X}\mathbf{v}
   = \lVert \mathbf{X}\mathbf{v}\rVert_2^2 \ge 0, \]
so the Hessian is positive semi-definite and the MSE is convex. A convex
function has no local minima --- only a global one --- so any stationary point
is \emph{the} minimum. Rearranging gives the \term{normal equation}:
\[ \hat{\boldsymbol\theta} \;=\;
   \left(\mathbf{X}^{\intercal}\mathbf{X}\right)^{-1}\mathbf{X}^{\intercal}y. \]
\end{derivbox}

Convexity is not a technicality to be waved through. It is also the reason
gradient descent, in Lecture 4, is guaranteed to reach the same answer this
closed form gives.

\subsection{What the stationarity condition says geometrically}

Read $\mathbf{X}^{\intercal}(\mathbf{X}\hat{\boldsymbol\theta} - y) =
\mathbf{0}$ again: every column of $\mathbf{X}$ is orthogonal to the residual
$\mathbf{r} = \mathbf{X}\hat{\boldsymbol\theta} - y$.

As $\boldsymbol\theta$ varies, $\mathbf{X}\boldsymbol\theta$ ranges over the
column space of $\mathbf{X}$. Minimising
$\lVert \mathbf{X}\boldsymbol\theta - y\rVert_2$ therefore finds the point of
that subspace closest to $y$ --- and the closest point of a subspace to a
vector is its orthogonal projection.

\begin{keybox}[title={The sentence to remember}]
Fitting a linear model is projecting the label vector onto the span of the
features. The residual is what the features cannot reach.
\end{keybox}

\subsection{Computed on our own data}

Eight numeric features, imputed and scaled, with a dummy column added, so
$\hat{\boldsymbol\theta}$ has nine entries. Solving the normal equation
directly and comparing with the library:

\begin{center}
\small
\begin{verbatim}
>>> theta_best[:3].round(1)
array([206333.5, -85371.4, -90829.7])
>>> lin_reg = LinearRegression().fit(X, y)
>>> np.abs(theta_best - np.r_[lin_reg.intercept_, lin_reg.coef_]).max()
6.4e-10
\end{verbatim}
\end{center}

Identical to floating-point precision. Look at $\theta_0$: \$206{,}334 --- and
the mean price of the training set is \$206{,}334. With centred features the
intercept \emph{is} the mean of the target, so the constant predictor we
measure later in this lecture is exactly this model with all eight weights set
to zero. Least squares spends those eight weights to get from \$115{,}311 down
to \$68{,}233.

\subsubsection{Verify the geometry, on the same data}

If the projection argument is right, the residual must be orthogonal to every
column of $\mathbf{X}$, so $\mathbf{X}^{\intercal}\mathbf{r}$ should be
numerically zero. It comes out at $7.1\times10^{-6}$ --- against
$\lvert\mathbf{X}^{\intercal}y\rvert_{\max} = 3.4\times10^{9}$, which is the
scale to judge it by. The relative size is $2\times10^{-15}$: machine epsilon
on a double.

\begin{keybox}[title={Two lessons in one check}]
The first is that this is a cheap, decisive test that your solver did what the
mathematics says. The second is more general: \emph{``small'' is meaningless
until you divide by something}. $7.1\times10^{-6}$ is enormous if the natural
scale is $10^{-9}$ and invisible if it is $10^{9}$.
\end{keybox}

\subsection{When it breaks}

The equation requires $(\mathbf{X}^{\intercal}\mathbf{X})^{-1}$ to exist.
$\mathbf{X}^{\intercal}\mathbf{X}$ is $(n{+}1)\times(n{+}1)$ and is invertible
if and only if it has full rank --- equivalently, if and only if $\mathbf{X}$
has linearly independent columns. There are exactly two ways to lose that.

\begin{failbox}[title={Failure condition --- more features than instances}]
If $m < n$ then $\operatorname{rank}(\mathbf{X}) \le m < n+1$, so
$\mathbf{X}^{\intercal}\mathbf{X}$ is singular. Note what the problem is: not
that no parameter vector fits the data, but that infinitely many fit it
\emph{exactly}. Common in genomics, in text, and in any wide-feature setting.
\end{failbox}

\begin{failbox}[title={Failure condition --- collinear features}]
If one column is a linear combination of the others, the columns are dependent
and $\mathbf{X}^{\intercal}\mathbf{X}$ is singular again.

This is Lecture 1's warning, now explained. ``When creating new combined
features, avoid simple weighted sums of existing features'' --- because a
weighted sum of existing columns \emph{is} a linear combination, and you would
be manufacturing the singularity yourself.
\end{failbox}

Exact collinearity is rare; near-collinearity is common, and is also a problem.
The matrix is invertible but ill-conditioned: tiny changes in the training set
produce large changes in $\hat{\boldsymbol\theta}$, and the model is
numerically unstable. Lecture 5 repairs this with ridge regression, which adds
$\alpha\mathbf{I}$ and makes the matrix invertible for \emph{every}
$\alpha > 0$ --- a fact that is immediate once you have this derivation.

\subsection{What Scikit-Learn actually does}

It does not invert anything. It computes $\hat{\boldsymbol\theta} =
\mathbf{X}^{+}y$, where $\mathbf{X}^{+}$ is the Moore--Penrose pseudoinverse,
obtained from the singular value decomposition:
\[ \mathbf{X} = \mathbf{U}\boldsymbol\Sigma\mathbf{V}^{\intercal}
   \qquad\Longrightarrow\qquad
   \mathbf{X}^{+} = \mathbf{V}\boldsymbol\Sigma^{+}\mathbf{U}^{\intercal}. \]
To build $\boldsymbol\Sigma^{+}$: take $\boldsymbol\Sigma$, set to zero every
value below a small threshold, replace the rest by their reciprocals, and
transpose.

The pseudoinverse is always defined, even when the inverse is not. But notice
what the threshold is doing. When the inverse does not exist there are
infinitely many $\boldsymbol\theta$ that fit equally well; zeroing the tiny
singular values selects one of them --- the minimum-norm least-squares
solution, the shortest vector among all the best fits.

\begin{keybox}[title={A choice, not a discovery}]
That selection is a convention, not a fact about the data. It is why two
rank-deficient fits can agree on every prediction and disagree on every
coefficient, and why interpreting coefficients from a collinear design is
unsafe even when the predictions are fine.
\end{keybox}

\begin{center}
\small
\begin{tabular}{lll}
\toprule
\textbf{Method} & \textbf{In $n$ (features)} & \textbf{In $m$ (instances)} \\
\midrule
Normal equation      & $O(n^{2.4})$ to $O(n^{3})$ & $O(m)$ \\
SVD (Scikit-Learn)   & $O(n^{2})$                 & $O(m)$ \\
\bottomrule
\end{tabular}
\end{center}

Both are linear in the number of instances, so both handle large training sets
well provided the data fits in memory. Both become very slow as $n$ grows ---
say 100{,}000 features. That is what gradient descent is for.

\section{Measuring generalisation}\label{sec:generalisation}

\subsection{A number that means nothing}

Fit an unconstrained decision tree on the 16{,}512 training districts and score
it on those same districts: RMSE $0.0$. No error at all. There are two
explanations and only one is flattering.

\emph{The model is perfect.} Housing prices would have to be an exact
deterministic function of nine census attributes, with two identical districts
always carrying identical median prices, no measurement noise anywhere in a
census, and no influence from anything outside our nine columns. Reject.

\emph{The model memorised.} An unconstrained tree keeps splitting until every
leaf is pure, and with enough depth it isolates every training district into
its own leaf. Predicting the training set is then a lookup, and a lookup has
zero error.

Such a model is called \term{nonparametric} --- not because it has no
parameters, but because their number is not fixed before training, so the
structure is free to follow the data. That freedom is exactly what lets it
memorise.

\begin{failbox}[title={Failure condition --- a score computed on the rows it was fitted to}]
Any score computed on the rows a model was fitted to measures how well it
reproduces what it has already seen. For a model flexible enough to memorise,
that is not evidence about anything else.

Linear regression on the same rows gives \$68{,}233, which happens to be honest
to within \$49 --- but only because the model is too rigid to memorise, and
\emph{nothing in the number tells you which case you are in}.
\end{failbox}

\subsection{``Then use the test set''}

No --- not yet. We are still choosing between models and tuning them. Every
time a choice is made by looking at the test score, that score stops being an
estimate of performance on unseen data and becomes a target we are fitting to.
The test set is used \emph{once}, at the end, to estimate the generalisation
error of a system already committed to.

This is the same argument as Lecture 1's specimen exam question, where a loop
chose $\alpha$ by scoring on \code{X\_test}. We answer the rest of it in
\S\ref{sec:specimen}.

\subsection{Leakage: the same error, one step earlier}

A score can be corrupted before any model is fitted.

\begin{center}
\small
\begin{verbatim}
scaler = StandardScaler()
housing_scaled = scaler.fit_transform(housing_full)   # all 20,640 rows
train_set, test_set = train_test_split(housing_scaled)  # then split
\end{verbatim}
\end{center}

\begin{definition}[Data leakage]
A quantity computed from rows outside the training set enters the training
path.
\end{definition}

Here the mean and standard deviation were computed from a set that includes the
test rows, so the model is evaluated on rows whose own values helped define the
transformation applied to them. It raises no exception, and the reported score
improves.

\subsubsection{What the leak cost here --- measured}

How much it improves is not knowable from the code. It is measurable
afterwards. Same data, same models, twenty seeds, scaled before the split
against scaled after it; all figures are dollars of RMSE.

\begin{center}
\small
\begin{tabular}{lrrr}
\toprule
\textbf{Model} & \textbf{seed 42} & \textbf{mean, 20 seeds} & \textbf{sd} \\
\midrule
Linear regression   &  0.88 & 0.02 &  0.31 \\
Random forest       & 57.17 & 3.89 & 47.87 \\
PCA(4), then linear & 15.90 & 1.85 & 16.34 \\
\bottomrule
\end{tabular}
\end{center}

Against a split-to-split spread of \$1{,}391 on the same RMSE, every one of
those is nothing --- and the sign flips from seed to seed. Report the seed-42
column alone and you have published noise.

\begin{keybox}[title={Why this is an argument \emph{for} the rule}]
The cost here was nothing. The cost elsewhere is unbounded. And you cannot tell
which case you are in by reading the code --- which is precisely why the rule
has to be procedural rather than a judgement made case by case.

Lecture 1 predicted this: all three models here are scale-invariant or nearly
so, and a standardisation is an invertible affine map. The leak was real and
its effect was zero, for a reason we could state in advance.
\end{keybox}

\begin{failbox}[title={The rule, stated once}]
Fit every transformer on the \emph{training set only}. Never call \code{fit()}
or \code{fit\_transform()} on the validation set, the test set, or new data.
Once fitted, you may \code{transform()} anything.

This applies to imputers, scalers, encoders, dimensionality reduction ---
everything that learns anything from data. We will not rely on remembering it;
the rest of this lecture builds an object that makes violating it structurally
impossible.
\end{failbox}

\subsection{What we need, then}

\begin{center}
\small
\begin{tabular}{p{0.46\textwidth}p{0.42\textwidth}}
\toprule
\textbf{Requirement} & \textbf{Because} \\
\midrule
An estimate of error on unseen rows & a training score cannot see overfitting \\
\ldots without touching the test set & it is spent the first time it influences
                                       a choice \\
\ldots with a \emph{spread}, not one number & two estimates differ only if the
                                       spread says so \\
\ldots and all preprocessing refitted inside it & otherwise the estimate leaks
                                       and is optimistic \\
\bottomrule
\end{tabular}
\end{center}

All four are satisfied by one procedure.

\section{Cross-validation, pipelines and tuning}

\subsection{From holdout to $k$-fold}

Hold out a validation set from the training set: train candidates on the
reduced set, select on the validation set, retrain the winner on the full
training set, evaluate once on the test set. The difficulty is sizing it. Too
small and the evaluations are imprecise, so you may select a poor model by
chance; too large and the remaining training set is much smaller than the one
the final model will use --- like selecting a sprinter by watching a marathon.

The resolution is many small validation sets, used in rotation. Split the
training set into $k$ non-overlapping \term{folds}, then train and evaluate $k$
times, holding out a different fold each time:
\[ \text{score} \;=\; \frac{1}{k}\sum_{j=1}^{k}
   \mathrm{RMSE}\!\left(\text{model trained without fold } j,\ \text{fold }
   j\right). \]
The cost is $k$ fits. The return is not only a mean but a \emph{spread}, which
a single holdout cannot give at any price.

\begin{keybox}[title={Two details in the call that are not decoration}]
\begin{verbatim}
cv = KFold(n_splits=10, shuffle=True, random_state=42)
tree_rmses = -cross_val_score(tree_reg, X_train, y_train,
                 scoring="neg_root_mean_squared_error", cv=cv)
\end{verbatim}
The minus sign: Scikit-Learn's cross-validation expects a utility (greater is
better), so the scorer is the \emph{negative} RMSE and the sign is flipped
back.

The explicit \code{KFold}: passing \code{cv=10} takes ten contiguous blocks in
row order. Two people whose frames are sorted differently then get different
folds and cannot see why their numbers disagree.
\end{keybox}

\subsection{The tree, measured honestly}

\begin{center}
\small
\begin{tabular}{lr}
\toprule
\textbf{Statistic over 10 folds} & \textbf{RMSE} \\
\midrule
mean   & 68{,}573.73 \\
min    & 65{,}499.05 \\
median & 69{,}070.93 \\
max    & 70{,}958.42 \\
\bottomrule
\end{tabular}
\end{center}

The honest number for that same tree is \$68{,}574, not zero. And the folds run
from \$65{,}499 to \$70{,}958, so any comparison turning on less than a couple
of thousand dollars is not a comparison.

\subsection{All three models, honestly}

\begin{center}
\small
\begin{tabular}{lrr}
\toprule
\textbf{Model} & \textbf{Training RMSE} & \textbf{Cross-validated RMSE} \\
\midrule
Predict a constant & \$115{,}311 & \$115{,}300 \\
Linear regression  & \$68{,}233  & \$68{,}282 \\
Decision tree      & \$0         & \$68{,}574 \\
Random forest      & \$18{,}058  & \$48{,}687 \\
\bottomrule
\end{tabular}
\end{center}

The gap between the columns is the overfitting, and the first row --- which
cannot overfit anything --- has no gap at all. Reading down:

\begin{itemize}[leftmargin=1.4em,itemsep=2pt]
  \item \term{Linear regression:} the columns nearly agree. It is not
        overfitting; it is \emph{under}fitting.
  \item \term{Decision tree:} \$0 against \$68{,}574. Severe overfitting ---
        and once measured honestly it is not better than the linear model it
        appeared to crush.
  \item \term{Random forest:} \$18{,}058 against \$48{,}687. Still
        overfitting, but the best honest score, and 58\% below the constant
        baseline.
\end{itemize}

\subsection{Is \$68{,}574 worse than \$68{,}282?}

The tree's mean is \$292 above the linear model's and the folds scatter by
\$2{,}000. But both models saw the \emph{same} ten folds, so compare them fold
by fold and the shared fold-to-fold noise cancels in the difference.

\begin{center}
\small
\begin{tabular}{lrrl}
\toprule
\textbf{Per-fold difference} & \textbf{Mean} & \textbf{Std} &
\textbf{Same sign in all 10?} \\
\midrule
Decision tree $-$ linear regression & $+\$292$     & \$1{,}634 & no \\
Random forest $-$ linear regression & $-\$19{,}594$& \$1{,}650 & yes \\
\bottomrule
\end{tabular}
\end{center}

The forest is genuinely better. The tree is \emph{indistinguishable} from the
linear model: a paired $t$-test gives $p = 0.59$, and on six of the ten folds
the tree is the better of the two.

\begin{failbox}[title={Failure condition --- the habit this is training}]
Say ``the tree is worse'' and you have read a ranking out of noise. Two numbers
are only different if you have shown that they are --- and the paired
comparison costs nothing, because you already ran both models on the same
folds.
\end{failbox}

Why does the forest win? It averages many trees fitted on random subsets, and
if their errors are not strongly correlated the average is more accurate than
any of them. Lecture 7 makes that precise by deriving the variance of an
average of correlated predictors.

\subsection{Pipelines make leakage structurally impossible}

\begin{center}
\small
\begin{verbatim}
full_pipeline = Pipeline([
    ("prep",  preprocessing),                       # the ColumnTransformer
    ("model", RandomForestRegressor(random_state=42)),
])
\end{verbatim}
\end{center}

When \code{cross\_val\_score} holds out fold $j$, the \emph{entire} pipeline
--- imputer, encoder, scaler, model --- is fitted on the other folds only. The
imputer's median and the scaler's mean never see the held-out fold. Leakage
cannot occur by accident, which is a stronger guarantee than any amount of
care.

\subsection{Tune, instead of guessing}

\begin{center}
\small
\begin{verbatim}
param_grid = {"model__max_features": [4, 6, 8, 10, 12],
              "model__n_estimators": [30, 100, 200]}
grid_search = GridSearchCV(full_pipeline, param_grid, cv=cv,
                 scoring="neg_root_mean_squared_error", n_jobs=-1)
\end{verbatim}
\end{center}

The double underscore addresses any hyperparameter of any estimator, however
deeply nested: \code{model\_\_max\_features} reads as ``the step named
\code{model}, its parameter \code{max\_features}''.

The winner is \code{max\_features=8}, \code{n\_estimators=200}, scoring
\$48{,}180 --- improved from \$48{,}687, a gain of \$508 against a fold spread
of \$2{,}762.

\begin{failbox}[title={Failure condition --- two things to notice, and neither is the improvement}]
150 fits bought a gain we cannot cleanly distinguish from noise.

And 200 was the \emph{largest} value tried for \code{n\_estimators}. A winner
sitting on the edge of the grid means the optimum may lie outside it: we did
not find the optimum, we found the edge of what we searched.
\end{failbox}

Preprocessing has hyperparameters too, and the same path reaches them:

\begin{center}
\small
\begin{tabular}{lr}
\toprule
\textbf{Imputation strategy} & \textbf{Cross-validated RMSE} \\
\midrule
\code{median}        & \$48{,}180 \\
\code{mean}          & \$48{,}154 \\
\code{most\_frequent}& \$48{,}171 \\
\bottomrule
\end{tabular}
\end{center}

A spread of \$26 across the three. With 207 missing values out of 20{,}640
that is exactly what it should be --- and now we know instead of assuming. The
reason to put preprocessing inside the pipeline is not that tuning it always
helps; it is that preprocessing and model interact, and the pipeline is what
lets you find out whether they do \emph{here}.

Randomised search is usually the better spend: it explores as many distinct
values as it has iterations rather than the few you listed, adding a
hyperparameter that turns out not to matter costs it nothing where a grid
multiplies, and --- note --- it cannot sit on the edge of a grid, because there
is no grid. Halving searches go further, running a tournament that evaluates
many candidates on a small slice and promotes the survivors with more data each
round.

\subsection{Analyse the best model}

\begin{center}
\small
\begin{tabular}{lr}
\toprule
\textbf{Feature} & \textbf{Importance} \\
\midrule
\code{median\_income}            & 0.4418 \\
\code{ocean\_proximity\_INLAND}  & 0.1511 \\
\code{longitude}                 & 0.1146 \\
\code{latitude}                  & 0.1043 \\
\code{housing\_median\_age}      & 0.0494 \\
\multicolumn{2}{c}{\ldots} \\
\code{ocean\_proximity\_NEAR OCEAN} & 0.0062 \\
\code{ocean\_proximity\_NEAR BAY}   & 0.0020 \\
\code{ocean\_proximity\_ISLAND}     & 0.0003 \\
\bottomrule
\end{tabular}
\end{center}

Median income dominates, as the correlation column predicted. Longitude and
latitude together take 0.22: the forest has rediscovered the coast from two raw
coordinates.

Importance is a reason to \emph{try} dropping a feature and then measure ---
never a reason to drop it outright, and never a measurement in itself.

\subsubsection{The gun loaded in Lecture 1}

At the bottom of that list: \code{ocean\_proximity\_ISLAND}, at 0.0003. Five
districts in California, and the split was stratified on income, not on this.

\begin{center}
\small
\begin{tabular}{lrrr}
\toprule
& \textbf{Total} & \textbf{In training} & \textbf{In test} \\
\midrule
\code{ISLAND} districts & 5 & 2 & 3 \\
\bottomrule
\end{tabular}
\end{center}

The model has seen the category twice. It will be asked about it three times,
and those three answers will be reported inside an average over 4{,}128.

\begin{failbox}[title={Failure condition --- and in cross-validation it can vanish entirely}]
\code{handle\_unknown="ignore"} does not raise on a category it never saw. It
emits $[0,0,0,0]$ --- a district that is neither near the ocean nor inland ---
silently.

With our seed no fold loses it. Over 500 reshufflings, 11.2\% produce at least
one fold in which the encoder is fitted without \code{ISLAND} and then asked
about it: a one-in-nine chance of a silently wrong column, invisible in the
score. This is what ``rare categories will cause trouble'' meant.
\end{failbox}

\section{The test set, once}

\begin{center}
\small
\begin{verbatim}
final_model = grid_search.best_estimator_   # refitted on all of X_train
final_predictions = final_model.predict(X_test)
final_rmse = root_mean_squared_error(y_test, final_predictions)   # $49,037
\end{verbatim}
\end{center}

We \code{predict} on the test set. We never \code{fit} anything on it.

\subsection{How precise is that estimate?}

A bootstrap over the squared errors gives a 95\% interval of
\$46{,}752--\$51{,}379, or roughly $\pm\$2{,}313$.

Why bootstrap rather than a $t$-interval? Not the skew --- it is 8.9, but with
$n = 4{,}128$ the central limit theorem covers the mean anyway. It is that we
report the \emph{square root} of the mean, and converting an interval for one
into an interval for the other needs the delta method. The bootstrap skips it.

The tuned cross-validated estimate was \$48{,}180 and the test result is
\$49{,}037: \$858 apart, against an interval of $\pm\$2{,}313$. They agree.
There is nothing here to explain.

\begin{failbox}[title={Failure condition --- the temptation, and why it is fatal}]
If you tuned a lot, test performance is often slightly worse than the
cross-validated estimate, because the system was fine-tuned to the validation
data. When that happens you must \emph{not} tweak hyperparameters to improve
the test number. Those improvements would not generalise; you would be
overfitting the test set, which is the only thing standing between you and a
confident fiction.

We searched 15 configurations, so the effect here is invisible. Search 10{,}000
and it will not be.
\end{failbox}

\subsection{Look at the ten worst predictions}

\begin{center}
\small
\begin{tabular}{rrrcl}
\toprule
\textbf{Actual} & \textbf{Predicted} & \textbf{Error} &
\textbf{income\_cat} & \textbf{ocean\_proximity} \\
\midrule
\$500{,}001 & \$129{,}715 & $-\$370{,}286$ & 2 & \code{<1H OCEAN} \\
\$500{,}001 & \$130{,}590 & $-\$369{,}412$ & 3 & \code{INLAND} \\
\$131{,}300 & \$464{,}117 & $+\$332{,}817$ & 5 & \code{<1H OCEAN} \\
\$500{,}001 & \$171{,}694 & $-\$328{,}307$ & 2 & \code{<1H OCEAN} \\
\$500{,}001 & \$175{,}534 & $-\$324{,}467$ & 1 & \code{INLAND} \\
\bottomrule
\end{tabular}
\end{center}

Seven of the ten are capped districts the model under-predicts by a quarter of
a million dollars. The third row is the mirror image: a district whose
\code{median\_income} is 15.0001 --- the \emph{other} cap --- and which the
model therefore prices as if it were Beverly Hills. It is worth \$131{,}300.

Both caps were spotted in a histogram before any model existed, and they are
now the two largest errors of the finished system.

\subsection{The average hides the model}

One number for 4{,}128 districts is a summary, not a finding.

\begin{center}
\small
\begin{tabular}{lrr@{\hskip 2em}lrr}
\toprule
\textbf{income\_cat} & \textbf{n} & \textbf{RMSE} &
\textbf{ocean\_proximity} & \textbf{n} & \textbf{RMSE} \\
\midrule
1 (poorest) &  165 & \$62{,}606 & \code{INLAND}     & 1{,}250 & \$39{,}829 \\
2           & 1{,}316 & \$43{,}082 & \code{<1H OCEAN}  & 1{,}862 & \$48{,}967 \\
3           & 1{,}447 & \$48{,}381 & \code{ISLAND}     &       3 & \$54{,}754 \\
4           &  728 & \$53{,}152 & \code{NEAR OCEAN} &   569 & \$57{,}191 \\
5 (richest) &  472 & \$54{,}333 & \code{NEAR BAY}   &   444 & \$60{,}195 \\
\bottomrule
\end{tabular}
\end{center}

The poorest 165 districts are predicted worst in dollars --- and their median
price is \$91{,}300, so in relative terms it is far worse still. \code{NEAR
BAY} is 51\% worse than \code{INLAND}. And the \code{ISLAND} number is computed
from three districts: report it with the count or do not report it.

A single RMSE of \$49{,}037 describes none of these groups. That is what
``check fairness'' means operationally, and it is a different report to the
client.

\subsection{Drop the capped districts?}

The cap causes the worst errors, so remove the 205 capped test districts and
measure again:

\begin{center}
\small
\begin{tabular}{lrrr}
\toprule
\textbf{Evaluated on} & \textbf{n} & \textbf{Model RMSE} &
\textbf{Constant baseline} \\
\midrule
All test districts     & 4{,}128 & \$49{,}037 & \$115{,}727 \\
Capped districts removed & 3{,}923 & \$44{,}197 & \$97{,}908 \\
\bottomrule
\end{tabular}
\end{center}

\begin{failbox}[title={Failure condition --- these two numbers cannot be compared}]
\$44{,}197 is not an improvement on \$49{,}037. We did not change the model ---
we changed the \emph{question}. The second row answers ``how well does it do on
districts below the cap?'', which is a different and easier question.

Notice that the baseline fell by 15\% as well, which is the tell. Any time a
number improves, ask first whether the evaluation set moved. Dropping rows is
the cheapest way in the world to improve a metric, and it is almost never a
result.
\end{failbox}

Should we drop them? For \emph{training}, arguably yes --- they teach the model
a price that does not exist. For \emph{reporting}, only with the sentence above
attached.

\subsection{The criterion we never measured}

The brief said the experts are off by more than 30\% \emph{of the price}. We
optimised dollars. Measured on the training folds --- because comparing two
differently-trained models is a choice, and choices are not made on the test
set:

\begin{center}
\small
\begin{tabular}{lrrr}
\toprule
\textbf{Trained on} & \textbf{RMSE} & \textbf{Median \% error} &
\textbf{Within 30\%} \\
\midrule
the price (what we did) & \$48{,}629 & 12.0\% & 85.0\% \\
the $\log$ of the price & \$49{,}646 & 11.5\% & 86.5\% \\
\bottomrule
\end{tabular}
\end{center}

Regressing the log target is \$1{,}017 worse in dollars and 1.5 points better
on the criterion the stakeholder actually stated. Neither is the right answer
--- but choosing without measuring is not a choice.

Also worth carrying to the client: the cheapest tenth of districts is predicted
with an RMSE of \$38{,}441, the dearest tenth with \$97{,}519. In dollars the
model is worst where the money is; in percentages it is worst where the money
is not.

\section{What the number means}

\begin{center}
\small
\begin{tabular}{lrl}
\toprule
\textbf{Measured how} & \textbf{RMSE} & \textbf{vs the baseline} \\
\midrule
Predict a constant --- the training mean & \$115{,}311 & --- \\
Best model, cross-validated on the training set & \$48{,}180 & 58\% better \\
The same model, on the test set, once & \$49{,}037 & 58\% better \\
\bottomrule
\end{tabular}
\end{center}

The first row is the one people leave out. Without it, \$49{,}037 is a number
with no units of judgement attached to it.

\subsection{Is that good?}

On the stakeholder's own criterion the system is within 30\% on 85.1\% of
districts, with a median miss of 11.3\%. The experts were said to be off by
more than 30\% typically. So: plausibly better --- on an assumption we were
handed and never checked.

\begin{keybox}[title={What a defensible report says}]
Ask for fifty districts estimated by the expert team and score them on the same
test rows with the same metric. Until then the comparison is between a
measurement and a claim, and only one of them has an interval.

It may still be worth launching, especially if it frees the experts for work
only they can do. But do not deploy it unqualified for the poorest districts,
or for the 205 at the cap. ``Is the model good?'' is not a question about RMSE.
It is a question about what the organisation does with it.
\end{keybox}

\subsection{Six questions to ask of any reported number}

\begin{enumerate}[leftmargin=1.6em,itemsep=2pt]
  \item Was every transformer fitted \emph{after} the split, on training data
        only?
  \item Is preprocessing inside the pipeline passed to cross-validation?
  \item Was the reported metric computed on held-out data?
  \item Was the test set touched more than once?
  \item Were hyperparameters selected using validation data, not test data?
  \item Does any feature encode information unavailable at prediction time?
\end{enumerate}

Questions 1 and 2 are what Part B of the paper asks when it hands you a code
excerpt.

\subsection{Delegating the review}

An assistant can perform step 3 of Lecture 1's loop, but only if you ask it as
an \emph{adversary} rather than an author. ``Is this code correct?'' invites
agreement, and you will usually get it. Instead: \emph{``This notebook reports
RMSE \$49{,}037. Assume that number is optimistic. List every path by which
test-set information could have reached the training procedure, and cite the
line.''} Give it the failure to look for and a number to disbelieve.

\begin{center}
\small
\begin{tabular}{lp{0.40\textwidth}}
\toprule
\textbf{Bug} & \textbf{Found from the code alone?} \\
\midrule
\code{fit\_transform} on the full set & yes \\
Preprocessing outside the cross-validated pipeline & yes \\
Test set evaluated more than once & no --- that is session history, not code \\
A feature unavailable at prediction time & no --- that is domain knowledge \\
Split not stratified on the right variable & no --- ``right'' is your judgement \\
\bottomrule
\end{tabular}
\end{center}

Three of the five require something no reader of the file can supply. Those
three are yours.

\subsection{Back to the specimen question}\label{sec:specimen}

Lecture 1 set a loop choosing $\alpha$ by scoring on \code{X\_test}. Parts 1
and 2 were answerable then; parts 3 and 4 are answerable now.

\emph{The repair.} Wrap the choice in a grid search over folds carved from the
training set, then touch the test set exactly once, after $\alpha$ is fixed.
\code{RidgeCV} does the same job in one line and is equally acceptable. What is
not acceptable is any version in which \code{X\_test} appears before the choice
is made.

\emph{Is the honest number higher?} In expectation, yes: the printed value was
a minimum over five noisy estimates, and $\alpha$ was chosen to look good on
the very set it was scored on. Guaranteed on one split? No --- the bias is a
statement about repeated draws, not about this sample, and cross-validation may
also pick a genuinely better $\alpha$. Both halves are required; ``higher''
alone, or ``no'' alone, is an incomplete answer.

You have now seen the second half measured: our own tuned model came out \$858
higher on the test set, well inside the interval --- a single split, telling
you nothing about the expectation.

\section{The notebook}

\begin{trybox}[title={What to change}]
The notebook repeats the Lecture 1 load and split in its first two cells, so it
stands on its own. Run it once end to end, then change the grid's
\code{n\_estimators} list to include values above 200 and re-run the search.
The winner moved off the edge of the grid --- did the score move with it, or by
less than the \$2{,}762 fold spread? That one change is the whole argument
about grid boundaries.
\end{trybox}

\section{Exercises}

Five, in the style of the written paper. Solutions on Lecture 3's deck.

\begin{enumerate}[leftmargin=1.6em,itemsep=4pt]
  \item The normal equation is
        $\hat\theta = (X^{\mathsf T}X)^{-1}X^{\mathsf T}y$. Your design matrix
        has a column equal to the sum of two others. State what happens to
        $X^{\mathsf T}X$, and what \code{np.linalg.pinv} returns instead.
        \hfill \scope{[5]}
  \item A pipeline scales the features and then fits a ridge model, and is
        passed whole to \code{cross\_val\_score}. Explain what would go wrong
        if the scaling were done once, before the call. \hfill \scope{[5]}
  \item You report a 95\% confidence interval for the test RMSE. A colleague
        says the interval proves the model is better than the baseline. Give
        the one question you would ask before agreeing. \hfill \scope{[4]}
  \item $k$-fold cross-validation reports a mean of 0.71 with folds
        $[0.52, 0.79, 0.75, 0.74, 0.75]$. What do you conclude, and what would
        you do next? \hfill \scope{[5]}
  \item Explain, in two sentences, why the test set may be looked at exactly
        once, and what you may legitimately do if the number disappoints.
        \hfill \scope{[4]}
\end{enumerate}

\section*{Summary}
\addcontentsline{toc}{section}{Summary}

Fitting a linear model is orthogonal projection: the minimiser satisfies
$\mathbf{X}^{\intercal}(\mathbf{X}\hat{\boldsymbol\theta}-y)=\mathbf{0}$, so
the residual is orthogonal to every feature, and
$\hat{\boldsymbol\theta} = (\mathbf{X}^{\intercal}\mathbf{X})^{-1}
\mathbf{X}^{\intercal}y$ whenever the inverse exists. It fails to exist when
$m < n$ or when the features are collinear --- which is why Lecture 1 warned
against building features as weighted sums of existing ones. The SVD
pseudoinverse is always defined, is cheaper, and silently chooses one of
infinitely many fits when the problem is degenerate.

A score computed on the training rows cannot detect overfitting: an
unconstrained tree scored $0$ and was worth \$68{,}574. The test set is spent
the first time it influences a choice. Leakage is a condition, not an accident,
and the pipeline makes it structurally impossible rather than merely
discouraged --- though measured here, over twenty seeds, the leak cost nothing,
for a reason Lecture 1 could state in advance.

$k$-fold gives an honest estimate \emph{and} its precision; two estimates
differ only if the paired per-fold difference says so, which demoted the tree
from ``much worse'' to ``indistinguishable, $p = 0.59$''. Tuning bought \$508
against a fold spread of \$2{,}762 and left the winner on the edge of the grid.
The test set was touched once, at \$49{,}037 with a 95\% interval of
$\pm\$2{,}313$, and then the average was taken apart: by income, by ocean
proximity, by the worst ten errors, and against a criterion the client stated
and we had not measured.

\end{document}
